% HƯỚNG DẪN SỬ DỤNG SÁCH
\chapter*{Hướng Dẫn Sử Dụng Sách}
\addcontentsline{toc}{chapter}{Hướng Dẫn Sử Dụng Sách}
\markboth{Hướng Dẫn Sử Dụng Sách}{Hướng Dẫn Sử Dụng Sách}

\begin{truyen}{Bộ Sách Này Dành Cho Ai?}{Who Is This Book For?}
\chuhoa{B}{ộ sách <<Truyện Tích Cảm Hứng Song Ngữ>> được thiết kế cho ba nhóm đối tượng chính.}
\textit{(The ``Inspiring Bilingual Tales'' series is designed for three main audiences.)}

\textbf{Giáo viên:} dùng mỗi câu chuyện ngắn để mở bài, kết tiết, hoặc kể vào những phút cuối giờ.
\textit{(\textbf{Teachers:} use each short story to open a lesson, close a class, or tell in the final minutes of a session.)}

\textbf{Học sinh:} đọc từng câu chuyện và dừng lại suy ngẫm --- không cần đọc liên tiếp, mỗi truyện là một đơn vị độc lập.
\textit{(\textbf{Students:} read each story and pause to reflect --- no need to read continuously, each story stands alone.)}

\textbf{Phụ huynh:} dùng như nội dung kể chuyện buổi tối, thay thế cho màn hình điện thoại, để nuôi dưỡng thói quen đọc và suy nghĩ.
\textit{(\textbf{Parents:} use as bedtime storytelling content, replacing screens, to nurture a habit of reading and reflection.)}
\end{truyen}

\vspace{6pt}

\begin{tcolorbox}[
  colback=trapgold!8,
  colframe=trapgold!70!black,
  coltitle=black,
  fonttitle=\bfseries,
  title={Giải Thích Các Ô Màu Trong Sách \quad\textit{\small(Color-Box Guide)}},
  arc=4pt, boxrule=1pt, breakable, enhanced
]

\smallskip
\textbf{\color{trapbrown}Ô Nâu~--- Câu Chuyện Chính (\textit{Main Story})}\\
Phần tường thuật song ngữ Việt -- Anh. Đọc cả hai để vừa hiểu nội dung vừa luyện tiếng Anh. Nếu bận, chỉ cần đọc tiếng Việt; nếu muốn luyện, che tiếng Việt và đọc tiếng Anh.\\
\textit{(Bilingual narrative in Vietnamese and English. Read both to understand content and practise English. If short on time, read only Vietnamese; to practise, cover the Vietnamese and read only English.)}

\medskip
\textbf{\color{trapgreen}Ô Xanh~--- Bài Học (\textit{Lesson Learned})}\\
Một hoặc hai câu tóm tắt bài học cốt lõi. Đây là phần nên ghi lại vào sổ tay.\\
\textit{(One or two sentences summarising the core lesson. This is the part worth noting in a journal.)}

\medskip
\textbf{\color{trapgreen!60}Ô Xanh Nhạt~--- Tiếng Anh Ngắn (\textit{Easy English})}\\
Hai đến ba câu tiếng Anh để nhớ nhanh. Đọc to và lặp lại~--- đây là luyện tập ngôn ngữ hiệu quả nhất.\\
\textit{(Two to three English sentences to remember quickly. Read aloud and repeat~--- this is the most effective language practice.)}

\medskip
\textbf{\color{trapgold!80!black}Ô Vàng~--- Tự Học Nhanh (\textit{Quick Self-Study})}\\
Ba câu hỏi suy ngẫm. Dành 2--3 phút trả lời trong đầu, hoặc viết ra. Không có câu trả lời đúng duy nhất~--- quan trọng là dừng lại và nghĩ.\\
\textit{(Three reflection questions. Spend 2--3 minutes answering mentally or in writing. There is no single right answer~--- what matters is pausing to think.)}

\smallskip
\end{tcolorbox}

\vspace{6pt}

\begin{baihoc}
\textbf{Gợi Ý Đọc~--- Cách đọc hiệu quả nhất:} 1~truyện mỗi ngày, đọc vào buổi sáng trước khi lên lớp hoặc buổi tối trước khi ngủ. Đừng cố đọc hết cả chương trong một lần.
\textit{(\textbf{Reading Tip~--- Most effective way:} 1~story per day, read in the morning before class or in the evening before sleep. Do not try to finish a whole chapter in one sitting.)}
\end{baihoc}

\vspace{4pt}

\begin{ghinhoanh}
\textbf{Reading tips:}

Read 1 story a day --- slow and steady wins.

Cover the Vietnamese and read only English to practise.

Write down one sentence that stays with you.
\end{ghinhoanh}

\ngancach
